\documentclass[11pt]{article}
\usepackage[utf8]{inputenc}
\usepackage[T1]{fontenc}
\usepackage{lmodern}
\usepackage{amsmath,amssymb,amsfonts}
\usepackage{graphicx}
\usepackage{booktabs}
\usepackage{hyperref}
\usepackage{geometry}
\geometry{margin=1in}

\title{Mathematical Foundations of the Imagized Language Model (ILM)}
\author{ImagizedLanguageModel Project}
\date{\today}

\begin{document}
\maketitle

\begin{abstract}
The Imagized Language Model (ILM) learns semantic representations directly from glyph images of words or characters, bypassing traditional tokenization. Using Alpaca-style question--answer (QA) data, ILM jointly trains: (i) a glyph encoder mapping 128\,\(\times\)\,128 images to a latent vector, (ii) a multi-channel product codebook that yields discrete, memory-like indices per token, and (iii) contrastive objectives linking glyphs, words, and sentence semantics. This note formalizes the data pipeline, model components, and training losses, and shows correspondence with the implemented training code.
\end{abstract}

\section{Overview}
Given a corpus of QA pairs, ILM optimizes three aligned representations:
\begin{itemize}
  \item Glyph encoder \( f_g : \mathbb{R}^{3\times H\times W} \to \mathbb{R}^d \) (with \(H=W=128\)).
  \item Product codebook embedding \( f_c : t \mapsto \mathbf{e}_t \in \mathbb{R}^d \) with multi-channel discrete indices.
  \item Contrastive and structural objectives: (glyph\,\(\leftrightarrow\)\,code) InfoNCE, (question\,\(\leftrightarrow\)\,answer) InfoNCE, and Gram alignment.
\end{itemize}

Supervision arises from self-alignment: glyph\,\(\leftrightarrow\)\,code embeddings and question\,\(\leftrightarrow\)\,answer semantics.

\section{Data and Tokenization}
For a QA pair \((Q_i, A_i)\), language-dependent tokenization yields
\[
  T_Q = [t_1^Q,\ldots,t_m^Q],\quad T_A = [t_1^A,\ldots,t_n^A].
\]
English uses regex wordpiece tokenization; Chinese uses per-character segmentation over CJK code points. Each token \(t\) is rendered as a glyph image \(x_t \in [0,1]^{3\times 128\times 128}\), via a language-specific renderer (dynamic tiling for EN; centered font glyph for ZH).

\section{Glyph Encoder}
A lightweight CNN \( f_g \) maps a glyph to a \(d\)-dimensional embedding, followed by L2 normalization:
\[
  \mathbf{g}_t = \frac{f_g(x_t)}{\lVert f_g(x_t)\rVert_2} \in \mathbb{R}^d.
\]
Adaptive pooling provides invariance to the exact tiling layout used for shorter words.

\section{Product Codebook Embedding}
We define a multi-channel codebook with \(C\) channels and \(K\) codes per channel. For channel \(c\), the code vectors are
\[
  E_c = [\mathbf{e}_{c,1},\ldots,\mathbf{e}_{c,K}] \in \mathbb{R}^{K\times (d/C)}, \qquad c=1,\ldots,C.
\]
Each token \(j\) has learnable assignment logits \(A_j \in \mathbb{R}^{C\times K}\). With temperature \(\tau\), channel-wise selection probabilities are
\[
  p_{j,c,k} = \frac{\exp(A_{j,c,k}/\tau)}{\sum\limits_{k'} \exp(A_{j,c,k'}/\tau)}.
\]
The token embedding is the concatenation of channel aggregates:
\[
  \mathbf{e}_j = \operatorname{concat}_{c=1}^{C} \; \sum_{k=1}^{K} p_{j,c,k} \, \mathbf{e}_{c,k} \in \mathbb{R}^d, \qquad \hat{\mathbf{e}}_j = \frac{\mathbf{e}_j}{\lVert \mathbf{e}_j \rVert_2}.
\]
As \(\tau \to 0\), the selection becomes hard: \(c_j^*(c) = \arg\max_k A_{j,c,k}\). The discrete tuple \(\text{addr}(j) = (c_j^*(1),\ldots,c_j^*(C))\) serves as a compact, memory-like index for token \(j\).

\section{Sentence Embeddings}
Sentence embeddings are averages of token embeddings:
\[
  \mathbf{s}_Q = \frac{1}{|T_Q|} \sum_{t\in T_Q} \mathbf{e}_t, \qquad \mathbf{s}_A = \frac{1}{|T_A|} \sum_{t\in T_A} \mathbf{e}_t,
\]
with normalized forms \(\hat{\mathbf{s}}_Q = \mathbf{s}_Q/\lVert\mathbf{s}_Q\rVert_2\) and \(\hat{\mathbf{s}}_A = \mathbf{s}_A/\lVert\mathbf{s}_A\rVert_2\).

\section{Objective Functions}
\subsection{Glyph--Code InfoNCE}
For \(U\) unique tokens in a mini-batch, with cosine similarities \(S_{uv} = \hat{\mathbf{g}}_u^\top \hat{\mathbf{e}}_v\), the InfoNCE objective is
\[
  \mathcal{L}_{\text{img}} = -\frac{1}{U} \sum_{u=1}^{U} \log \frac{\exp(S_{uu}/\tau_g)}{\sum_{v=1}^{U} \exp(S_{uv}/\tau_g)}.
\]

\subsection{QA InfoNCE (Semantic Alignment)}
For \(B\) QA pairs, define similarities \(M_{ij} = \hat{\mathbf{s}}_{Q_i}^\top \hat{\mathbf{s}}_{A_j}\). The QA InfoNCE loss is
\[
  \mathcal{L}_{\text{qa}} = -\frac{1}{B} \sum_{i=1}^{B} \log \frac{\exp(M_{ii}/\tau_s)}{\sum_{j=1}^{B} \exp(M_{ij}/\tau_s)}.
\]

\subsection{Gram Matrix Structural Loss}
Using token-level normalized embeddings inside sentences, form
\(
  Z_Q = [\hat{\mathbf{e}}_{t_1^Q},\ldots,\hat{\mathbf{e}}_{t_m^Q}]^\top,\quad
  Z_A = [\hat{\mathbf{e}}_{t_1^A},\ldots,\hat{\mathbf{e}}_{t_n^A}]^\top.
\)
The Gram matrices are \(G_Q = Z_Q Z_Q^\top\) and \(G_A = Z_A Z_A^\top\). With \(m_0=\min(m,n)\),
\[
  \mathcal{L}_{\text{gram}} = \frac{1}{B} \sum_{i=1}^{B} \frac{1}{m_{0,i}^2} \big\lVert G_{Q_i}^{(m_{0,i})} - G_{A_i}^{(m_{0,i})} \big\rVert_F^2.
\]

\subsection{Code Usage Entropy Regularization}
Let \(P = \operatorname{softmax}(A, \text{dim}=K) \in \mathbb{R}^{V\times C\times K}\) be assignment probabilities over the vocabulary \(V\). The average entropy promotes diverse code usage:
\[
  H = -\frac{1}{VC} \sum_{j=1}^{V} \sum_{c=1}^{C} \sum_{k=1}^{K} P_{j,c,k} \log P_{j,c,k}, \qquad \mathcal{L}_{\text{reg}} = -\lambda_u H.
\]

\subsection{Total Loss}
Combining the terms (with typical hyperparameters \(\tau_g=\tau_s=0.07\), \(\lambda_g=0.1\), \(\lambda_u=0.01\)):
\[
  \boxed{\mathcal{L} = \mathcal{L}_{\text{img}} + \mathcal{L}_{\text{qa}} + \lambda_g\,\mathcal{L}_{\text{gram}} + \mathcal{L}_{\text{reg}}.}
\]

\section{Gradient Behavior and Learning Dynamics}
\begin{center}
\begin{tabular}{@{}lll@{}}
\toprule
Term & Gradient target & Semantic effect \\
\midrule
\(\mathcal{L}_{\text{img}}\) & glyph encoder \(f_g\), codebooks & visual--symbolic consistency \\
\(\mathcal{L}_{\text{qa}}\) & sentence means \(\mathbf{s}_Q,\mathbf{s}_A\) & semantic grouping/alignment \\
\(\mathcal{L}_{\text{gram}}\) & token geometry within Q/A & structural/grammatical consistency \\
\(\mathcal{L}_{\text{reg}}\) & assignment logits \(A\) & anti-collapse, code diversity \\
\bottomrule
\end{tabular}
\end{center}

\section{Consistency with Implementation}
The training step mirrors the formulas:
\begin{itemize}
  \item Stage 1: compute normalized glyph and code embeddings for unique tokens; apply InfoNCE on the \(U\times U\) cosine matrix.
  \item Stage 2: compute sentence means for QA pairs (normalized); apply QA InfoNCE on the \(B\times B\) cosine matrix.
  \item Gram alignment: build normalized token matrices for each Q and A; minimize Frobenius distance between Gram matrices (size matched by \(m_0\)).
  \item Regularization: add negative mean assignment entropy to encourage broad code usage per channel.
\end{itemize}

\section{Semantic Emergence and Discrete Memory Indices}
The product codebook induces a factorized latent space: each channel captures a complementary aspect (e.g., lexical category, domain, morphological pattern), and the concatenation composes them. InfoNCE terms pull together tokens that serve similar roles across QA, while Gram alignment encourages relational structure. The learned hard indices \(\text{addr}(j)\) then serve as a compact memory address per token, suitable for table lookups or hierarchical indexing.

\section{Diffusion Extension (Optional)}
To extend from static embeddings to generative modeling over glyph grids, one can train a UNet to invert a corruption process on glyph frames. With standard variance-preserving diffusion,
\[
  \mathbf{x}_t = \sqrt{\bar{\alpha}_t}\,\mathbf{x}_0 + \sqrt{1-\bar{\alpha}_t}\,\boldsymbol{\epsilon}, \quad \boldsymbol{\epsilon} \sim \mathcal{N}(0, I),
\]
the training objective is
\[
  \mathcal{L}_{\text{diff}} = \mathbb{E}_{t,\boldsymbol{\epsilon}}\,\big\lVert \boldsymbol{\epsilon} - \boldsymbol{\epsilon}_\theta(\mathbf{x}_t, t) \big\rVert_2^2,
\]
applied to glyph feature maps or latent embeddings. This yields coherent in-filling and iterative refinement of text-like images.

\section{Conclusion}
ILM jointly optimizes visual, structural, and semantic coherence by connecting glyph encodings with a discretized, factorized codebook and QA-driven alignment. The resulting representation learns semantics without external lexicons and produces discrete memory indices for efficient downstream use. The objectives above match the implemented pipeline.

\end{document}

